\section{Dataset and Refactoring Benchmark}

\subsection{Dataset A Overview}

Dataset~A is the core artifact of this paper: a maintenance-centred benchmark for nearby-version proof breakage in verified Dafny programs. Each benchmark item is organized around a verified seed program, a semantics-preserving refactored variant, the resulting verification outcome, and the outcomes of the baseline-recovery procedures introduced in Section~4. This definition makes the benchmark unit explicitly revision-aware. Rather than treating verification as a one-shot proving task, Dataset~A captures what happens when a previously verified program is subjected to controlled source evolution and then replayed under the same verification workflow. In that respect, it complements recent Dafny benchmark efforts aimed at benchmark generation, synthesis, or compositional reasoning rather than maintenance under evolution~\cite{DafnyBench2024,Xu2026DafnyComp,Wang2026TACoDafny}. The benchmark's value is therefore not only that it contains transformed programs, but that it preserves enough structure to support causal interpretation of how and why verification breaks.

\subsection{Seed Corpus and Strata}

The seed corpus is restricted to already verified Dafny programs whose proof-relevant structure is visible at the source level through contracts, assertions, loop invariants, and ghost code~\cite{Leino2010Dafny,Grov2016VerificationPatternsDafny}. The goal is not to maximize raw syntax diversity, but to assemble examples that expose plausible maintenance stress when their structure changes. To support later stratified analysis, Dataset~A groups seeds into three operational strata: tutorial and educational examples, benchmark or example-suite programs, and proof-engineering-heavy programs with substantial manual annotation. This split distinguishes lightweight pedagogical cases from instances that exercise more realistic proof maintenance behavior. It also keeps the benchmark from collapsing into a single style of Dafny development and makes it possible to ask whether nearby-version proof breakage concentrates in one source regime or persists across multiple levels of proof engineering maturity. Where available, ecosystem tooling and testing work provide additional evidence that Dafny now has a broad enough surrounding infrastructure to support this style of empirical study~\cite{Irfan2022TestingDafny,Fedchin2023ToolkitTestingDafny}.

\subsection{Benchmark Design}

The main design choice behind Dataset~A is to use controlled refactorings instead of unconstrained repository history. This choice trades away some messiness of real-world evolution in exchange for reproducibility, clearer causal interpretation, and benchmark items with cleaner ground truth. In the software-engineering literature, refactoring denotes structural change intended to preserve behavior~\cite{Mens2004RefactoringSurvey,Kim2012RefactoringInTheWild,Tsantalis2014RefactoringOperations}. This paper adopts that idea in a narrower form suited to verification maintenance: each benchmark item is built to isolate whether source-level structural change destabilizes an existing verified program even when the intended behavior is unchanged.

This makes Dataset~A scientifically different from generation-oriented Dafny benchmarks~\cite{DafnyBench2024,Xu2026DafnyComp,Wang2026TACoDafny}. Those benchmarks focus on constructing or solving verification tasks from scratch; Dataset~A instead measures maintenance under evolution. Its explicit axes are the seed-program stratum, the refactoring class applied to the seed, the post-refactoring verification outcome, and the baseline-recovery outcome recorded after replay. The benchmark therefore excludes deliberate semantic changes, verifier bugs as the primary object of study, and unconstrained edit histories whose causal structure is too noisy for the paper's central question.

\subsection{Refactoring Catalog}

P1 studies seven semantics-preserving refactoring classes chosen to span a range of plausible maintenance stresses rather than to exhaust all possible evolution events. The first group consists of local structural rewrites that mostly preserve control flow while perturbing proof-relevant names or local context: renaming changes the symbolic surface seen by transferred hints, helper extraction moves intermediate reasoning into a new abstraction boundary, and inlining collapses an existing abstraction back into its caller. A second group reshapes the local control or data context in which proofs discharge: legal statement reordering alters the order in which facts become available, non-behavioral loop reshaping perturbs loop-local reasoning without changing intended behavior, and contract factoring redistributes specification material across helper boundaries or shared clauses. The final group targets proof-context displacement directly: ghost-code movement relocates proof hints, witness structure, or auxiliary lemmas while leaving executable behavior unchanged.

These classes are retained because each plausibly perturbs a different proof-relevant mechanism while remaining in scope as semantics-preserving refactoring. Renaming and ghost-code movement primarily test surface continuity of proof hints. Helper extraction, inlining, and contract factoring stress how verification obligations depend on local abstraction boundaries. Legal statement reordering and loop reshaping stress sensitivity to source order, invariant placement, and nearby solver context. Collectively, the catalog is intended to cover a meaningful spread of maintenance pressures on verified Dafny programs without conflating them with semantic edits or underspecified programs. That breadth matters for interpretation in Section~5: if brittleness appears across these distinct categories, the result is harder to dismiss as a quirk of a single transformation style.

\subsection{Generation Pipeline}

Each benchmark item is produced by a replayable four-stage pipeline. First, a candidate seed is validated to ensure that the original Dafny program verifies successfully under the recorded toolchain. Second, an eligible controlled transformation is applied to create a semantics-preserving nearby-version variant. Third, the transformed program is re-verified and its outcome is captured, including success, failure, timeout, or exclusion when the transformation does not apply cleanly. Fourth, the resulting pair is packaged with benchmark metadata so that all baselines operate over the same original program, transformed program, and verification context.

This pipeline makes the artifact boundary explicit. For every retained item, Dataset~A stores the original program, the transformed program, the observed verification result after refactoring, a slot for a proof-oriented failure label, and slots for the baseline-recovery outcomes reported later in the paper. Cases that cannot be transformed into a valid nearby-version pair are logged and removed rather than silently repaired by hand. That policy keeps the benchmark interpretable as an evaluation substrate for maintenance under evolution and preserves comparability across baseline runs.

\subsection{Failure-Mode Labels and Benchmark Outputs}

Dataset~A is intended to support more than a single breakage percentage, so broken instances are annotated with a compact set of failure-mode labels. The label set used in P1 distinguishes broken assertion flow, weakened loop invariants, brittle lemma or helper-application structure, trigger or quantifier sensitivity, ghost-code or proof-hint displacement, and unexplained or tool-sensitive failure. These labels are analytic categories for empirical characterization rather than definitive mechanistic diagnoses. Their purpose is to support stratified reporting, to make the residual repair gap interpretable, and to indicate which broken cases are likely to require stronger forms of repair than simple replay or surface transfer.

At the artifact level, the benchmark outputs comprise the transformed instances themselves, their verification outcomes, their failure-mode labels when verification breaks, and the baseline-recovery outcomes and cost hooks used in Section~5. This makes Dataset~A reusable as both a dataset and an evaluation substrate: a future repair system can be evaluated on the same nearby-version pairs, while this paper focuses on establishing the benchmark composition and the residual gap left by simple baselines. Section~4 builds directly on this structure by defining how baseline recovery is measured over each benchmark item.
